\documentclass[11pt]{article}
\usepackage[a4paper,margin=1in]{geometry}
\usepackage{fontspec}
\usepackage[boldfont=false]{xeCJK}
\setCJKmainfont{FandolSong-Regular.otf}
\usepackage{amsmath}
\usepackage{amssymb}
\usepackage{array}
\usepackage{booktabs}
\usepackage{graphicx}
\usepackage{adjustbox}
\usepackage{enumitem}
\setlist[itemize]{topsep=2pt,partopsep=0pt,parsep=0pt,itemsep=2pt,leftmargin=1.4em}
\setlist[enumerate]{topsep=2pt,partopsep=0pt,parsep=0pt,itemsep=2pt,leftmargin=1.6em}
\usepackage{parskip}
\usepackage{fvextra}
\fvset{breaklines=true,breakanywhere=true,fontsize=\small}
\setlength{\parindent}{0pt}

\begin{document}

\begin{center}
  {\LARGE\bfseries Microeconomic decision-makers}\\[0.35em]
  {\large Igcse Economics}
\end{center}
\vspace{0.6em}

\section*{Money and banking}

Long ago, people swapped goods directly. This is called \textbf{barter} 物物交换. Barter is hard: you must find someone who has what you want \emph{and} wants what you have. \textbf{Money} 货币 solves this problem.

\subsection*{The functions of money}

Money does four jobs:

\begin{itemize}
\item a \textbf{medium of exchange} 交换媒介 — you swap money for goods, so you do not need barter.

\item a \textbf{store of value} 价值储藏 — you can hold money now and spend it later, because it keeps its value.

\item a \textbf{unit of account} 记账单位 — money measures and compares the value of things (this is the price).

\item a \textbf{means of deferred payment} 延期支付手段 — you can buy now and pay later.

\end{itemize}

\subsection*{The characteristics of money}

Good money must be:

\begin{itemize}
\item \textbf{durable} 耐用的 — it lasts a long time.

\item \textbf{portable} 便于携带的 — easy to carry.

\item \textbf{divisible} 可分割的 — can be split into small amounts for small payments.

\item limited in supply, hard to copy, and accepted by everyone.

\end{itemize}

\subsection*{Banks}

\begin{itemize}
\item A \textbf{central bank} 中央银行 is the government's bank. It issues notes, looks after the banking system, and sets \textbf{interest rates} 利率. There is one in each country.

\item \textbf{commercial banks} 商业银行 are the ordinary banks that people and \textbf{firms} 企业 use. They keep savings safe, lend money, and let people pay each other.

\end{itemize}

\section*{Households}

A \textbf{household} 家庭 is a person or group living together and making money decisions. Each household must choose how much to spend, save, and borrow.

\begin{itemize}
\item \textbf{saving} 储蓄 is income that is not spent now.

\item \textbf{borrowing} 借贷 is spending money you do not have yet, and paying it back later (usually with an extra charge).

\end{itemize}

Three main things change these choices:

\begin{itemize}
\item \textbf{income} 收入 — with more income, people spend and save more.

\item interest rates — a higher interest rate makes saving more rewarding and borrowing dearer. So saving rises and borrowing falls.

\item \textbf{confidence} 信心 — if people feel sure about their jobs and the future, they spend and borrow more. If they are worried, they save more.

\end{itemize}

\section*{Workers}

\subsection*{Choosing a job}

When a worker chooses an \textbf{occupation} 职业 (a type of job), two kinds of factor matter:

\begin{itemize}
\item \textbf{wage factors} — the pay. The \textbf{wage} 工资 is the money you earn. Higher pay attracts more workers.

\item \textbf{non-wage factors} — everything else: working hours, holidays, how safe or pleasant the job is, whether you enjoy the work, distance from home, and the chance to move up to a higher job.

\end{itemize}

\subsection*{Why workers earn different amounts}

Workers are not all paid the same. The reasons include:

\begin{itemize}
\item \textbf{skill} 技能 and \textbf{qualifications} 资格 — workers who are hard to replace, or who studied for years, earn more.

\item danger — risky or dirty jobs may pay more to attract workers.

\item \textbf{trade unions} 工会 — a trade union is a group of workers who join together to bargain for better pay and conditions. Strong unions can raise pay.

\item \textbf{discrimination} 歧视 — unfair treatment because of gender, race, or age can cause unfair pay gaps.

\end{itemize}

\subsection*{Division of labour}

\textbf{Division of labour} 劳动分工 (also called \textbf{specialisation} 专业化) means splitting a job into small tasks, with each worker doing one task again and again.

\textbf{Advantages}: workers get faster and better at their one task; less time is wasted moving between tasks; this raises output.

\textbf{Disadvantages}: doing one task is boring, so workers may make mistakes or leave; if one worker is absent, the whole line can stop; each worker learns only one skill.

\section*{Firms}

\subsection*{The size of firms}

Firms come in many sizes, from one person to huge companies. We can measure size by the number of workers, the amount of \textbf{capital} 资本 used, or the output.

Some firms stay small because: the \textbf{market} 市场 is small; the product is personal (like a hairdresser); the owner wants to stay in control; or there is little money to grow.

\subsection*{How firms grow}

Firms grow by selling more, or by joining with other firms. When two firms join, it is a \textbf{merger} 合并 (or takeover). There are three types:

\begin{itemize}
\item \textbf{horizontal merger} 横向合并 — two firms at the \emph{same} stage of the \emph{same} industry join (two car makers).

\item \textbf{vertical merger} 纵向合并 — two firms at \emph{different} stages of the same industry join (a car maker and a tyre maker). "Backward" is towards the supplier; "forward" is towards the seller.

\item \textbf{conglomerate merger} 混合合并 — two firms in \emph{different} industries join (a car maker and a food company).

\end{itemize}

\subsection*{Business failure}

Firms can suffer \textbf{business failure} 企业倒闭 (they close down) for many reasons: costs too high, too few sales, poor management, too many rivals, not enough cash, or a fall in \textbf{demand} 需求.

\section*{Firms and production}

\subsection*{Labour-intensive and capital-intensive}

\begin{itemize}
\item \textbf{labour-intensive} 劳动密集型 production uses a lot of \textbf{labour} 劳动 (many workers) and little machinery — for example, hand-made crafts.

\item \textbf{capital-intensive} 资本密集型 production uses a lot of capital (machines) and few workers — for example, a modern car factory.

\end{itemize}

\subsection*{Productivity}

\textbf{Productivity} 生产率 means output per worker (or per machine) in a set time. Higher productivity means each worker makes more, so the cost of each unit falls. Training, better machines, and good management all raise productivity.

\subsection*{Demand for and supply of factors}

Firms buy \textbf{factors of production} 生产要素 (\textbf{land} 土地, labour, capital, and \textbf{enterprise} 企业家才能). The price and amount of each factor are set by its own demand and supply, just like any market. For example, if many firms want skilled workers, the wage for those workers rises.

\subsection*{Economies and diseconomies of scale}

As a firm grows and makes more, the cost of each unit can change:

\begin{itemize}
\item \textbf{economies of scale} 规模经济 — as output rises, the average cost of each unit \emph{falls}. Reasons: buying materials in bulk is cheaper, big machines work more efficiently, and large firms borrow money cheaply.

\item \textbf{diseconomies of scale} 规模不经济 — if a firm grows \emph{too} big, the average cost of each unit starts to \emph{rise}. Reasons: a huge firm is hard to manage, messages travel slowly, and workers may feel less important.

\end{itemize}

These can be \emph{internal} (caused by the firm's own size) or \emph{external} (caused by the whole industry growing).

\section*{Firms' costs, revenue and objectives}

\subsection*{Costs}

A \textbf{cost} 成本 is money a firm spends to produce. Costs split into:

\begin{itemize}
\item \textbf{fixed costs} 固定成本 — costs that do \emph{not} change with output, such as rent and insurance. You pay them even if you make nothing.

\item \textbf{variable costs} 可变成本 — costs that \emph{change} with output, such as materials and power. More output means more variable cost.

\item \textbf{total cost} 总成本 — fixed cost plus variable cost.

\item \textbf{average cost} 平均成本 — the cost of \emph{each} unit: total cost divided by the number of units made.

\end{itemize}

\subsection*{Revenue}

\textbf{Revenue} 收益 is the money a firm gets from selling its goods.

\begin{itemize}
\item \textbf{total revenue} 总收益 — the price multiplied by the quantity sold.

\item \textbf{average revenue} 平均收益 — total revenue divided by the quantity sold (this is the same as the price).

\end{itemize}

\subsection*{Profit and other objectives}

\textbf{Profit} 利润 is what is left after costs:

profit = total revenue − total cost

Most firms want as much profit as possible. But firms can have other \textbf{objectives} 目标:

\begin{itemize}
\item \textbf{survival} 生存 — a new firm first just wants to stay open.

\item \textbf{growth} 增长 — getting bigger, to gain economies of scale and market power.

\item social objectives — some firms aim to help people or protect the environment, not only to earn profit.

\end{itemize}

\section*{Types of markets}

Markets differ by how many firms there are and how hard it is for new firms to join.

\subsection*{Competitive markets}

A competitive market has many firms selling similar goods, and it is easy for new firms to enter. This is strong \textbf{competition} 竞争. In such a market:

\begin{itemize}
\item firms must keep prices low and quality high, or buyers go elsewhere.

\item firms work hard to cut costs and make new products.

\end{itemize}

This is good for consumers: low prices, wide choice, and good quality.

\subsection*{Monopoly}

A \textbf{monopoly} 垄断 is a market with only \emph{one} firm, or one firm that controls most of it. New firms find it hard to enter because of \textbf{barriers to entry} 进入壁垒 — for example, very high start-up costs, or the one firm controlling key supplies.

In a monopoly:

\begin{itemize}
\item the firm can charge high prices, because buyers have no other choice.

\item there is little pressure to improve quality or cut costs.

\item but a large monopoly may gain economies of scale, and may use its profit to pay for research.

\end{itemize}

\subsection*{Competition versus monopoly}

\begin{center}
\adjustbox{max width=\linewidth}{%
\begin{tabular}{lll}
\toprule
\textbf{} & \textbf{Consumers} & \textbf{Firms} \\
\midrule
competition & lower prices, more choice, better quality & lower profit; must work hard to survive \\
monopoly & higher prices, less choice & higher profit and an easy life, but may grow lazy \\
\bottomrule
\end{tabular}}
\end{center}


\end{document}
